% !TeX root = ../../金融经济学_王江_讲义及习题.tex
% 来源：教材印刷页223—225，PDF物理页238—240；保留题号与数据。
% 原题14.1、14.3的不相关条件限于n不等于n'；14.3以列向量b_n和转置统一维度，变换后载荷为b'_n。
% 14.2公式恢复教材15章之前的原编号14.9—14.11，修正旧解漏写a_C；14.1有限N保留剩余方差项。

\Needspace{10\baselineskip}
\par\noindent\begin{minipage}{\linewidth}
\begin{problem}
\WJFourteenExercise{exer:wj-fourteen-1}{} 市场上有无穷多的证券。它们的收益率由单因子模型给定：
\[
    \tilde{r}_{n}=\bar{r}_{n}+b_{n}\tilde{f}+\tilde{u}_{n},\quad n=1,2,\ldots
\]
这里 $\mathrm{E}[\tilde{f}]=0$；$\forall\,n,n'$，$\mathrm{E}[\tilde{u}_{n}]=0$，$\mathrm{E}[\tilde{f}\tilde{u}_{n}]=0$；$\operatorname{Var}[\tilde{u}_{n}]\le\sigma^{2}$；以及 $\mathrm{E}[\tilde{u}_{n}\tilde{u}_{n'}]=0$（$n\ne n'$）。一个组合由 $\{w_{1},w_{2},\ldots\}$ 给定。如当 $N$ 很大时，
\[
    \sum_{n}^{N}w_{n}^{2}\simeq O(N^{-1})
\]
则称它是充分分散的，其中 $N$ 代表组合中绝对权重最大的股票数目而 $O(\cdot)$ 表示同阶或更高阶。
\begin{enumerate}[label=(\alph*)]
    \item 计算充分分散组合的方差。
    \item 证明它不受单个证券的特殊风险的影响。
    \item 更一般地讨论什么因素决定了充分分散组合的风险，这里风险的度量是收益率的方差。
\end{enumerate}
\end{problem}
\end{minipage}\par
\Needspace{5\baselineskip}
\begin{solution}

\begin{enumerate}[label=(\alph*)]
\item 记前 $N$ 只证券的组合载荷为 $b_{p,N}=\sum_{n=1}^Nw_nb_n$，剩余支付为 $\widetilde u_{p,N}=\sum_{n=1}^Nw_n\widetilde u_n$。由不相关条件，
\[
 \operatorname{Var}(\widetilde r_{p,N})
 =b_{p,N}^2\sigma_f^2+
 \underbrace{\sum_{n=1}^Nw_n^2\operatorname{Var}(\widetilde u_n)}_{\text{剩余方差}},
 \qquad 0\le\operatorname{Var}(\widetilde u_{p,N})\le\sigma^2\sum_{n=1}^Nw_n^2=O(N^{-1}).
\]

因此 $\operatorname{Var}(\widetilde r_{p,N})=b_{p,N}^2\sigma_f^2+O(N^{-1})$。
\item 剩余支付均值为零，且 $\mathrm E[\widetilde u_{p,N}^2]\to0$，所以它在均方意义下趋于零。充分分散消除的是组合剩余风险的极限影响。
\item 因子方差和组合载荷共同决定保留的系统风险。若 $b_{p,N}\to b_p$，则组合方差趋于 $b_p^2\sigma_f^2$；若组合载荷也趋于零，则总方差趋于零。\qedhere
\end{enumerate}
\end{solution}

\Needspace{10\baselineskip}
\par\noindent\begin{minipage}{\linewidth}
\begin{problem}
\WJFourteenExercise{exer:wj-fourteen-2}{} 假设资产 $A$、$B$、$C$ 的收益服从精确两因子模型：
% 仅习题模式省略正文公式，保留教材公式起始号；完整讲义接续正文自动编号。
\ifdefined\WJExercisesOnly\setcounter{equation}{8}\fi
\begin{align}
    r_{A}&=a_{A}+b_{A,1}\tilde{f}_{1}+b_{A,2}\tilde{f}_{2}\label{eq:wj-14-9}\\
    r_{B}&=a_{B}+b_{B,1}\tilde{f}_{1}+b_{B,2}\tilde{f}_{2}\label{eq:wj-14-10}\\
    r_{C}&=a_{C}+b_{C,1}\tilde{f}_{1}+b_{C,2}\tilde{f}_{2}\label{eq:wj-14-11}
\end{align}
其中 $\tilde{f}_{1}$ 和 $\tilde{f}_{2}$ 是两个均值为 0、方差为 1 的因子。$r_{F}$ 是无风险利率。为消除套利机会，$a$、$b$ 以及 $r_{F}$ 应该满足什么条件？
\end{problem}
\end{minipage}\par
\Needspace{5\baselineskip}
\begin{solution}

令 $a=(a_A,a_B,a_C)^{\mathsf T}$，$B$ 的三行为各资产的两因子载荷，$\iota$ 为全1列向量，$R=1+r_F$。无套利要求存在因子溢价 $\lambda=(\lambda_1,\lambda_2)^{\mathsf T}$，满足
\[
 \boxed{a-r_F\iota=B\lambda},\qquad
 a_i-r_F=b_{i,1}\lambda_1+b_{i,2}\lambda_2,\quad i=A,B,C.
\]

\textbf{推导：}取风险资产投资金额向量 $z$，并用无风险证券融资。其零成本终期支付为
\[
 z^{\mathsf T}(a-r_F\iota)+(B^{\mathsf T}z)^{\mathsf T}\widetilde f.
\]

若 $B^{\mathsf T}z=0$，支付为常数；无套利要求 $z^{\mathsf T}(a-r_F\iota)=0$。因此 $a-r_F\iota\in(\ker B^{\mathsf T})^\perp=\operatorname{col}B$，即上式。

也可取权重 $(x_1,x_2,1-x_1-x_2)$ 的成本为1的组合。其期望收益为
\[
 a_C+(a_A-a_C)x_1+(a_B-a_C)x_2.
\]

当两因子载荷均消为零时，该期望收益应等于 $r_F$。这与上面的矩阵条件一致，并适用于载荷矩阵秩不足的情形。

因子收益分布还决定带有因子风险的头寸是否可能形成套利。若 $\widetilde f$ 在 $\mathbb R^2$ 上有完全支持，上述条件也充分：非零因子载荷的支付有正、负两种可能，零因子载荷的零成本支付为零。一般分布下，完整无套利条件可表为存在严格正定价核 $m$，满足
\[
 \mathrm E[m]=R^{-1},\qquad
 B\,\mathrm E[m\widetilde f]=-R^{-1}(a-r_F\iota).
\]

它同时约束预期收益、因子载荷和因子支付的支持集。\qedhere
\end{solution}

\Needspace{10\baselineskip}
\par\noindent\begin{minipage}{\linewidth}
\begin{problem}
\WJFourteenExercise{exer:wj-fourteen-3}{} 考虑一个有无穷多资产和 $K$ 个因子的因子模型：
\[
    \tilde{r}_{n}=\bar{r}_{n}+b_{n}^{\mathsf T}\tilde{f}+\tilde{u}_{n},\quad n=1,2,\ldots
\]
其中 $\tilde{f}=\begin{bmatrix}\tilde{f}_{1} \\ \ldots \\ \tilde{f}_{K}\end{bmatrix}$ 是因子向量，$b_{n}=\begin{bmatrix}b_{n,1} \\ \ldots \\ b_{n,K}\end{bmatrix}$ 是资产 $n$ 的因子载荷向量，而 $\forall\,n,n'$，$\mathrm{E}[\tilde{f}]=0$，$\mathrm{E}[\tilde{u}_{n}]=0$，$\mathrm{E}[\tilde{f}\tilde{f}^{\mathrm{T}}]=I$，$\mathrm{E}[\tilde{f}\tilde{u}_{n}]=0$，$\mathrm{E}[\tilde{u}_{n}\tilde{u}_{n'}]=0$（$n\ne n'$），其中 $I$ 是单位阵。考虑另外一组因子 $\tilde{f}'$，它们是 $\tilde{f}$ 的线性变换：
\[
    \tilde{f}'=\mathrm{A}\tilde{f}
\]
这里 $\mathrm{A}$ 是一个 $(K\times K)$ 的矩阵。
\begin{enumerate}[label=(\alph*)]
    \item 证明存在 $\tilde{f}'$ 使得
    \[
        \tilde{r}_{n}=\bar{r}_{n}+b'_{n}{}^{\mathsf T}\tilde{f}'+\tilde{u}_{n},\quad n=1,2,\ldots
    \]
    其中 $\mathrm{E}[\tilde{f}']=0$，$\mathrm{E}[\tilde{f}'\tilde{f}'^{\mathrm{T}}]=I$，$\mathrm{E}[\tilde{f}'\tilde{u}_{n}]=0$，$\forall\,n$。因此，因子模型不是唯一的。
    \item 证明因子模型有许多因子不同但却等价的表述形式。
\end{enumerate}
\end{problem}
\end{minipage}\par
\Needspace{5\baselineskip}
\begin{solution}

\begin{enumerate}[label=(\alph*)]
\item 任取正交阵 $A$，令 $\widetilde f'=A\widetilde f$、$b'_n=Ab_n$。则
\[
 b'_n{}^{\mathsf T}\widetilde f'
 =b_n^{\mathsf T}A^{\mathsf T}A\widetilde f=b_n^{\mathsf T}\widetilde f,
\]
\[
 \mathrm E[\widetilde f']=0,\quad
 \mathrm E[\widetilde f'\widetilde f'{}^{\mathsf T}]=AA^{\mathsf T}=I,\quad
 \mathrm E[\widetilde f'\widetilde u_n]=A\mathrm E[\widetilde f\widetilde u_n]=0.
\]

因子、载荷同时旋转，收益支付及其协方差保持不变。
\item 任取可逆阵 $T$，令 $\widetilde g=T\widetilde f$、$d_n=T^{-\mathsf T}b_n$，仍有 $d_n^{\mathsf T}\widetilde g=b_n^{\mathsf T}\widetilde f$；此时因子协方差为 $TT^{\mathsf T}$。其中正交变换还保留单位协方差。$K\ge2$ 时有无穷多个正交旋转；$K=1$ 时标准化表述有正、负两个方向，一般可逆变换还允许非零尺度变化。\qedhere
\end{enumerate}
\end{solution}

\Needspace{10\baselineskip}
\par\noindent\begin{minipage}{\linewidth}
\begin{problem}
\WJFourteenExercise{exer:wj-fourteen-4}{*} 给定练习~\ref{exer:wj-fourteen-3} 描述的因子模型，APT 要求
\[
    \bar{r}_{n}-r_{F}=\lambda^{\mathrm{T}}b_{n}
\]
其中 $\lambda=\begin{bmatrix}\lambda_{1} \\ \lambda_{2} \\ \ldots \\ \lambda_{K}\end{bmatrix}$ 是因子风险溢价向量。证明存在一个等价的因子模型，它的因子向量为 $\tilde{f}^{*}$ 而对应的因子风险溢价向量为 $\lambda^{*}=\begin{bmatrix}\lambda_{1}^{*} \\ 0 \\ \ldots \\ 0\end{bmatrix}$，即只有第一个因子有风险溢价而其它因子虽然代表了不可分散风险，却没有风险溢价。
\end{problem}
\end{minipage}\par
\Needspace{5\baselineskip}
\begin{solution}

若 $\lambda=0$，取 $\widetilde f^*=\widetilde f$、$b_n^*=b_n$ 即可。以下设 $\lambda\ne0$，令 $q_1=\lambda/\|\lambda\|$，将其补成一组正交单位基，组成正交阵 $Q=(q_1,\ldots,q_K)$。

定义 $\widetilde f^*=Q^{\mathsf T}\widetilde f$、$b_n^*=Q^{\mathsf T}b_n$、$\lambda^*=Q^{\mathsf T}\lambda$，则
\[
 b_n^{*\mathsf T}\widetilde f^*=b_n^{\mathsf T}\widetilde f,\qquad
 \lambda^{*\mathsf T}b_n^*=\lambda^{\mathsf T}b_n,
 \qquad\boxed{\lambda^*=(\|\lambda\|,0,\ldots,0)^{\mathsf T}}.
\]

正交变换保留单位因子协方差及与剩余项的不相关性。新坐标的第一个方向与风险溢价向量平行，其余方向与之正交，因此风险价格为零；这些方向的因子方差仍为1。共同风险的大小与其价格由此分别体现。\qedhere
\end{solution}

